\textbf{COLUNI 2018 -- Questão 35}\par\medskip Em 25 de agosto de 1961, o Presidente Jânio Quadros renunciou à presidência da República. Esse fato causou uma crise política, pois alguns setores das Forças Armadas e políticos da União Democrática Nacional (UDN) queriam impedir a posse do vice-presidente João Goulart na presidência da República, conforme estava previsto na Constituição de 1946.

\par

A solução política encontrada pelo Congresso Nacional para resolver essa crise foi:\begin{enumerate}
\item Apoiar a intervenção das Forças Armadas através de um golpe militar impedindo a posse de João Goulart.
\item Apoiar um Movimento Popular de Mobilização Nacional conhecido como Campanha da Legalidade.
\item Aprovar uma Emenda Constitucional mudando o Sistema Presidencialista para o Sistema Parlamentarista.
\item Convocar novas eleições presidenciais para legitimar a eleição direta de João Goulart.\end{enumerate}
